no claim about how the factor fork is normalized and installs no additional
inner witness.  That strictly later task belongs to the factor/return closure
of Section~\ref{sec:P2closed}; if it eventually needs its first inner token,
Lemma~\ref{lem:seed-policy} governs the one possible use of $\zeta$.

Suppose instead that the obligation continues canonically.  With
$y=A(x)$ and $g=1-\alpha(x)$, Lemma~\ref{lem:phase-aligned-side} gives
\[
 b=\chi_g(y).
\]
Thus the next obligation is aligned with the token just installed.
Lemma~\ref{lem:aligned-catch} says that within this next macro either its
ordinary source becomes exactly $b$, or $b$ is strictly replaced before the
route becomes attached.  In particular there is no tokenless multi-$A$
streak after the first seed.

Once an ordinary obligation source itself is a retained finite token,
Lemma~\ref{lem:finite-front} applies to every $A$-selecting continuation and
forces a strict token replacement (or an attached exceptional constructor)
before any later high return.  Hence an arbitrarily long canonical $A$-streak
is a finite concatenation of strict token events and attached routing; it is
not an equal-rank walk through the numerical cursors
$x,A(x),A^2(x),\ldots$.  This is exactly the point at which the old
pre-streak estimate failed.

It remains to inspect a valuation-two $B$ return from a retained finite
source $q$.  Corollary~\ref{cor:exceptional-no-b2} excludes all four
exceptional residue classes.  Put $y=B(q)$ and let the transferred phase be
$g=\alpha(q)$.  If $q$ is LOSS, then $y$ is a WIN child with
$h(y)\le h(q)-1$, so $q\mapsto y$ is strict.  If $q$ is WIN, choose a
canonical LOSS child.  When that child is $y$, the same strict replacement
holds.  In the remaining orientation the canonical LOSS child is $A(q)$.
The ordinary side identity gives
\[
 c=B(A(q))\in\operatorname{moves}(A(q))\cap\operatorname{moves}(y),
 \qquad h(c)\le h(q)-2.
\]
Lemma~\ref{lem:second-selector} says that the transferred phase at $y$
selects this exact $c$.  Therefore the reset is entered already carrying the
strict descendant $c$; it is either caught by the aligned/source-token case
or treated as an attached task.  No temporary cursor is promoted.

Finally, if a canonical streak reaches a $B$-selecting phase only after
several $A$ steps, the preceding argument is still exhaustive: immediately
after the first outer seed the next obligation is aligned; within at most one
canonical step its source catches the retained token unless that token has
already decreased.  Every subsequent $A$-obligation whose source is finite
spends a descendant by Lemma~\ref{lem:finite-front}.  Thus every
outcome-compatible streak length is covered without a numerical bound on the
pre-streak source.  The mandatory regression
$20\to30\to45\to68$, $B(68)=25>20$, is therefore harmless: the proof never
uses the comparison $25<20$ (which is false).
\end{proof}


\section{Semantic closure of factor, return and attached routing}
\label{sec:P2closed}

The remaining task is to prove that an attached token cannot disappear when
the arithmetic passes through a factor fork, a high return, a marked tail or
a short terminal lift.  We state the local facts in the form needed for the
global rank.  For audit provenance, the same arithmetic identities also occur with longer
case-by-case derivations in Sections 91--135 of the frozen proof ledger
\texttt{docs/verified-results.md} at repository revision
\nolinkurl{9d12789de3dbd841851f3faff1a237018d43879d}.  The proofs below reproduce
the incidences and rank payments actually used, so that ledger is a
cross-check rather than a logical black box.  Sections 137--138 of the ledger
are not used here; the audited entry argument below replaces them.

\begin{lemma}[Finite-source export]\label{lem:finite-export}
Let $O(x,\alpha(x))$ be an $A$-selecting obligation, let its common side be
\WIN, and suppose the ordinary source $x$ is finite.
\begin{enumerate}[label=(\alph*)]
\item If $x$ is \LOSS, the canonical branch exports an actual \LOSS token
      at least two proof levels below $x$ before any later high return.
\item If $x$ is \WIN and either $x$ is ordinary or $A(x)$ is the
      canonical \LOSS child, the return carries a \WIN token of height at
      most $h(x)-2$.
\item In the sole remaining orientation, $x$ is exceptional and $B(x)$ is
      its canonical \LOSS child; the returned adjacent pair is an exact
      factor-free lift over that token, whose height is $h(x)-1$.
\end{enumerate}
\end{lemma}

\begin{proof}
Put $y=A(x)$ and let $b$ be the common \WIN side of the obligation.  If
$x$ is \LOSS, then $y$ is \WIN with $h(y)\le h(x)-1$.  Since one child of
$y$ is the \WIN state $b$, its other child is \LOSS and has height at most
$h(x)-2$.

If $x$ is \WIN, choose its canonical \LOSS child $\ell_0$.  When $x$ is
ordinary, Lemma~\ref{lem:side} makes $p=B(A(x))$ a child of $\ell_0$ in the
orientation where $\ell_0=B(x)$; if $\ell_0=A(x)$, $p$ is directly a child
of that LOSS state.  In either case $p$ is \WIN and
$h(p)\le h(x)-2$.  When $x$ is exceptional and
$\ell_0=B(x)$, Lemma~\ref{lem:exceptional-side} gives the two returned
children as adjacent lifts over $\ell_0$.  These are exactly the three rows.
\end{proof}

\begin{lemma}[Generic lower factor forks spend the retained mark]
\label{lem:factor-fork-pay}
Consider the factor alternative of a typed $A$-selecting obligation, with the
finite provenance supplied by that alternative.  Then every lower factor
exponent pays or remains attached before control can restart.

For lower exponent $r=1$, let $b$ be the WIN factor source, let $V$ be the
retained LOSS state from the factor alternative and put $D=B(V)$.  If $y$ is
the ordinary child of $b$ selected by the factor phase, the exact twin identity
is
\[
 \{D,y\}=\{A(b),B(b)\},\qquad D\ne y.
\]
Hence $D$ is WIN, $y$ is the actual LOSS child of $b$, and the row either
exposes a boundary WIN at least two proof levels below $b$ or transfers the
DRAW to an adjacent frame over this same LOSS token.

For every $r\ge2$, let $x$ be the finite ordinary source from which the
factor fork was generated.  Then $x$ is ordinary and the factor common side is
$b=B(A(x))$, a child of both $A(x)$ and $B(x)$.  Splitting the actual outcome
of $x$ gives respectively: a strict canonical coefficient-source edge for the
same DRAW state; a common WIN child of a canonical LOSS child with height at
most $h(x)-2$; or the LOSS state $A^2(x)$ with height at most $h(x)-2$.
\end{lemma}

\begin{proof}
In the $r=1$ factor alternative the standard twin calculation gives
\[
 \{D,y\}=\{A(b),B(b)\},\qquad D\ne y.
\]
Since $V$ is LOSS, its
child $D=B(V)$ is WIN.  The factor source $b$ is also WIN, so its other child
$y$ is necessarily LOSS and $h(b)=h(y)+1$.  If the factor phase is
A-selecting, the common side of the exponent-one/two frame is a child of $y$
and therefore is WIN with height at most $h(b)-2$.  A continuing DRAW member
of that frame points to this lower WIN boundary.  If the phase is B-selecting,
Lemma~\ref{lem:b-diamond} transfers the DRAW to the exact adjacent frame over
the selected source $y$, retaining the actual LOSS token.  No unproved
finite outcome is introduced.

For $r\ge2$, the signed transition at $A(x)$ is B-selecting, so the selected
source is $b=B(A(x))$.  The four exceptional residue classes would make this
phase A-selecting with valuation one; hence this row is ordinary.  By
Lemma~\ref{lem:side}, $b$ is a child of both ordinary children of $x$.

If $x$ is DRAW, no new game state is inferred: the \emph{same actual DRAW
state} $x$, written in its canonical constant-tail coordinates, has coefficient
source $\rho(x)<x$ by Lemma~\ref{lem:source-bound}.  Replacing the retained
inner numerical anchor $x$ by its genuine canonical source is therefore a
strict source edge.

If $x$ is WIN, choose a canonical LOSS child of $x$.  The common state $b$ is
a child of it, hence is WIN and $h(b)\le h(x)-2$.  If $x$ is LOSS, both
ordinary children are WIN; the factor row already makes $b$ WIN, so the other
child $A^2(x)$ of $A(x)$ is LOSS and has height at most $h(x)-2$.  These are
all outcome orientations and every finite-token comparison is a genuine
parent/child comparison.
\end{proof}

\begin{lemma}[High-return token descent]\label{lem:high-return-pay}
Assume the high-return guard, including the exact obligation
$O(c,\alpha(c))$ and the finite return provenance supplied by the preceding
factor scan. A high return can never be a pure reset.
\begin{enumerate}[label=(\alph*)]
\item In each short return row $v=5,6$, the exact finite return token
      $c=Q_{v-3}^g(3J(t))$ is already marked with its proved finite outcome.
      The common-side split either replaces a retained token by a proper
      descendant or enters the exact exceptional lift over its canonical
      lower LOSS token.
\item For every $v\ge7$ the return guard supplies finite WIN states
\[
 u=Q_{v-1}^g(J(t)),\qquad c=Q_{v-3}^g(3J(t)),
\]
with $A(u)$ LOSS. Put $p=B(A(u))=A(c)$ and let $b=B(p)$. Then
\[
 h(p)\le h(u)-2,\qquad h(b)\le h(c)-2.
\]
Thus the pair replacement $(u,c)\mapsto(p,b)$ is strict before either
continuation of $O(c,\alpha(c))$ is followed.

If the obligation takes its canonical continuation, the next ordinary
source is exactly $p=A(c)$, so the strict token $p$ is inherited and no
marked-tail task is installed. If it takes its factor alternative, there are
further finite states $q,y$ such that
\[
 (p,b)\longmapsto(q,y)
\]
is componentwise strict, the continuing DRAW lies in the exact frame
\[
 \{Q_3^g(J(b)),Q_2^g(J(b))\},
\]
and, with $D=v-6$ and $C=27J(t)$,
\[
 b=Q_D^g(C),\qquad
\begin{array}{c|cc}
 v&q\text{ WIN}&y\text{ LOSS}\\ \hline
 7&B(b)&A(b)\\
 8&A(b)&B(b)\\
 v\ge9&A(b)&B(b).
\end{array}
\]
Hence only the factor alternative enters the marked-tail control, and it
does so after a strict pair-token replacement.
\end{enumerate}
\end{lemma}

\begin{proof}
For $v=6$ the common side is $b=B(A(c))$. If this side is DRAW, the
constant-tail state $c$ cannot be WIN: its canonical LOSS child would also
have $b$ as a child and would force $b$ WIN. Hence $c$ is LOSS. The WIN child
$A(c)$ has the DRAW child $b$, so its other child $A^2(c)$ is LOSS with
$h(A^2(c))\le h(c)-2$.

For $v=5$ the common side is $b=A^2(c)$. If $b$ is DRAW and $c$ is LOSS,
then $A(c)$ is WIN and its other child $B(A(c))$ is LOSS at least two levels
below $c$. If $c$ is WIN, then $A(c)$ is not LOSS because it has the DRAW
child $b$; therefore $B(c)$ is the canonical LOSS child of $c$. In the
ordinary row $B(A(c))$ is a WIN child of this LOSS token, while in the
exceptional row Lemma~\ref{lem:exceptional-side} gives the exact adjacent
lift over it. If the common side is WIN instead, Lemma~\ref{lem:finite-export}
applies because the statement explicitly retains the finite outcome of $c$.
Thus both short rows pay or remain attached to an already lower token.

Now let $v\ge7$. Direct constant-tail reduction gives
\[
 p=B(A(u))=A(c).
\]
The return guard makes $A(u)$ LOSS, so $p$ is its WIN child and
$h(p)\le h(u)-2$. The state $c$ is ordinary in this range. Since $p=A(c)$
is WIN, the other child $B(c)$ is LOSS, and
\[
 b=B(p)=B(A(c))
\]
is a child of that LOSS state. Hence $b$ is WIN and
$h(b)\le h(c)-2$. This proves the first strict pair replacement before the
obligation branch is selected.

In the canonical alternative of $O(c,\alpha(c))$, the next source is
$A(c)=p$ exactly. Thus the already lower token is inherited and the route
returns directly to an obligation over $p$.

In the factor alternative, $p$ is ordinary. Because $p$ and its child $b$
are WIN, the other child of $p$ is LOSS. Its ordinary common child
$q=B(A(p))$ is WIN and satisfies $h(q)\le h(p)-2$. The factor phase selects
the other child $y$ of the WIN state $b$; since $q$ is already its WIN child,
$y$ is the unique LOSS child and $h(y)=h(b)-1$. Hence
$(p,b)\mapsto(q,y)$ is strict in the multiset order.

Finally direct substitution gives $b=Q_{v-6}^g(27J(t))$. The signed
transition at $b$ selects $A(b)$ at $v=7$, $B(b)$ with valuation at least
three at $v=8$, and $B(b)$ with valuation two for every $v\ge9$. In the
factor alternative the selected state is exactly the LOSS token $y$ and the
other child is the WIN token $q$, yielding the displayed orientation table.
Thus the factor branch, and only that branch, enters the marked-tail frame
with an actual continuing DRAW and with the required tokens already marked.
\end{proof}

\begin{lemma}[Every long marked tail pays its LOSS token]
\label{lem:marked-tail-pay}
Let $C$ be positive and odd, let $D\ge3$, and put
\[
 b=Q_D^g(C),\qquad q=A(b),\qquad y=B(b).
\]
Suppose the marked high-return constructor retains the exact outcomes and
DRAW frame
\[
 b,q\text{ \WIN},\qquad y\text{ \LOSS},\qquad
 \{Q_3^g(J(b)),Q_2^g(J(b))\}\text{ contains a \DRAW}.
\]
Put
\[
 r=\begin{cases}
    B(q)=B(y),&D=3,\\
    B(q)=A(y),&D\ge4.
   \end{cases}
\]
Then $r$ is a \WIN child of the retained LOSS state $y$, so
\[
 h(r)\le h(y)-1,
\]
and the complete marked tail may replace $y$ by $r$ before any inner reset.
\end{lemma}

\begin{proof}
For $D=3$, the exponent-two/one common-child identity gives
$B(q)=B(y)$.  For $D\ge4$, Lemma~\ref{lem:long-tail} gives
$B(q)=A(y)$.  In either row $r$ is a child of the LOSS state $y$, hence is
WIN and has strictly smaller height.

The attachment survives the complete two-level scan by exact coordinates,
not by a comparison of canonical coefficients.  At $D=3$ the first raw
member is $Q_M^g(J(r))$ for its actual suffix exponent $M\ge2$.  For every
$D\ge4$ it is exactly
\[
 S=Q_1^g(J(r)).
\]
If the signed branch is taken, its common side is an exact lift over the
phase-selected ordinary child of $r$; if neither first exit is DRAW, the
second bridge is the obligation over this same $S$ and its transferred
source is again an exact lift over that selected child.  At the three short
endpoints $D=3,4,5$ these states coincide with the first/second signed and
raw exits of the original factor scan over $r$.  Hence no branch introduces
a source unrelated to $r$.

Therefore the strict replacement $y\mapsto r$ may be performed before the
raw/signed choice, after which all later routing data are allowed to reset:
the multiset rank has already decreased.  The displayed identities prove the exact attachment and use no source
comparison on a factorful cursor.  They agree with the marked-tail identities
of the frozen local ledger.
\end{proof}

\begin{lemma}[The two shortest marked tails]
\label{lem:short-tail}
Let $C$ be positive and odd, put $b=Q_D^g(C)$, and assume the exact marked
frame contains a continuing DRAW.

For $D=1$ the retained orientation is
\[
 q=B(b)\text{ WIN},\qquad y=A(b)\text{ LOSS}.
\]
If
\[
 9J(b)+1-2g=2^jJ(w),
\]
then $j\ge2$; moreover $w=A(y)$ when $j=2$ and $w=B(y)$ when
$j\ge3$. Hence the $D=1$ row always replaces the LOSS token $y$ by an
actual WIN child before any inner reset.

For $D=2$ the retained orientation is
\[
 q=A(b)\text{ WIN},\qquad y=B(b)\text{ LOSS}.
\]
If $y=0$, the returned arithmetic has coefficient source zero and is routed
by Lemma~\ref{lem:zero-source}.  If $y>0$, the first factor numerator has
valuation exactly one and its returned source has the exact form
\[
 w=Q_m^\delta(J(y)),\qquad \delta=1-g,\qquad m\ge2.
\]
Thus the row is not a free restart: it is a constant-tail lift over the
already retained LOSS token $y$.
\end{lemma}

\begin{proof}
For $D=1$, $b=2C-g$. Direct substitution gives
\[
9J(b)+1-2g=
\begin{cases}
2(27C+5),&g=0,\\
2(27C-5),&g=1.
\end{cases}
\]
Since $C$ is odd, the parenthesized integer is even, so $j\ge2$.  Moreover,
with $y=A(b)=3C-g$ the same identity is exactly
\[
 9J(b)+1-2g
 =2\bigl(3J(y)+1-2(1-g)\bigr).
\]
The expression in parentheses is the ordinary signed source transition at
$y$ in phase $1-g$.  If its valuation is one (total $j=2$) it selects
$A(y)$; if its valuation is at least two (total $j\ge3$) it selects $B(y)$.
Thus $w=A(y)$ for $j=2$ and $w=B(y)$ for $j\ge3$. Both are actual children
of the retained LOSS state $y$ and are therefore WIN proper descendants. In
particular the apparent $j=1$ lift row sometimes listed in a coarse
trichotomy is arithmetically empty here.

For $D=2$, $b=4C-g$ and
\[
9J(b)+1-2g=
\begin{cases}
2(54C+5),&g=0,\\
2(54C-5),&g=1.
\end{cases}
\]
The parenthesized integer is odd, so the valuation is exactly one.  Let
$z=3C-g$; by the exponent-two/one boundary identity $y=R(z)$.  The returned
source $w$ is determined by
\[
 J(w)=54C+5-10g,
\]
so $w=18C+1$ for $g=0$ and $w=18C-2$ for $g=1$.  Its canonical odd
coefficient is
\[
 \kappa(w)=\begin{cases}
 \odd(9C+1),&g=0,\\
 \odd(9C-1),&g=1.
 \end{cases}
\]
Lemma~\ref{lem:alt-lift}, applied to $3C$ in the first row and to $3C-1$ in
the second, gives $\kappa(w)=J(y)$, so no factor of three is hidden in the
coefficient of $w$.

When $y>0$, let $\lambda$ be the maximal alternating-suffix length of $z$.
For $g=0$ the alternating-lift identity is
$3z+1=2^\lambda J(y)$, whence
\[
 w=2(3z+1)-1=Q_{\lambda+1}^{1}(J(y)).
\]
For $g=1$ the same identity reads $3z+2=2^\lambda J(y)$, whence
\[
 w=2(3z+2)=Q_{\lambda+1}^{0}(J(y)).
\]
